\documentclass{article}

\usepackage{Mathematics}
\pdftitle{Block-E4-PDP2}

\title{\textbf{Block E4 \\ PDP2, HSLU, Semester 4}}
\author{Matteo Frongillo}
\date{}

\begin{document}

\part*{Block E4 - Effective prompting}
\section*{Good promts}
\subsection*{Rules to follow}
\begin{itemize}
    \item Start with the simplest request
    \item Use clear separators between different tasks, such as
        \begin{verbatim}
            ### INTRODUCTION ###
            ### OUTPUT STYLE ###
        \end{verbatim}
    \item Use clear words to tell the AI what you want, for instance ``Translate'', ``List'', ``Explain''
    \item Be very specific as you add details: ``Explain the differece between renewable and non-renewable energy to a 16-year-old. Use simple words. Include three examples each.''
    \item Go to the point and remove all the unnecessary informations.
    \item Do not be generic and specify what you want: ``\sout{Write few paragraphs}''.\\
        Better: ``Write 2 paragraphs'' or ``Write a 200 words essay''
    \item Avoid to say what not to do, instead say what to do.
\end{itemize}

\subsection*{Good prompt example}
\begin{verbatim}
You are a librarian

### RULES ###
Before starting generating the output, ask:
- How old are you?
- What is your favorite book?
- How many book per year do you read?
- What is your favorite book genre?

### GOAL ###
Suggest a book based on the user age, favorite book, number of books read the past year,
and book genre.

### OUTPUT ###
- Return a list of 5 books
- Add a 30 words description for each book
- Ask the user if they want 5 more books suggested
- Do not ask further questions

### CONSTRAINTS ###
- Do not suggest the same book
- Do not suggest books from the same saga, if any
\end{verbatim}

\subsection*{Bad prompt example}
\begin{verbatim}
Suggest me some books Tolkien-style.
\end{verbatim}


\end{document}